\newcommand{\todo}[1]{{\color{red}{#1}}}

\newcommand{\mk}{\textsc{miniKanren}\xspace}
\newcommand{\micro}{\textsc{microKanren}\xspace}
\newcommand{\merc}{\textsc{Mercury}\xspace}
\newcommand{\curry}{\textsc{curry}\xspace}
\newcommand{\haskell}{\textsc{Haskell}\xspace}
\newcommand{\ocaml}{\textsc{OCaml}\xspace}
\newcommand{\ocanren}{\textsc{OCanren}\xspace}
\newcommand{\prolog}{\textsc{Prolog}\xspace}
\newcommand{\inm}{\emph{in}\xspace}
\newcommand{\outm}{\emph{out}\xspace}

\newcommand{\LLang}{\mathcal{L}}
\newcommand{\semantics}[2]{\llbracket #1 \rrbracket_{#2}}
\newcommand{\LSem}[1]{\semantics{#1}{\LLang}}

\newcommand{\code}[1]{\texttt{#1}}

\definecolor{lstKeyword}{HTML}{000080}
\definecolor{lstBG}{HTML}{FFFFFF}
\definecolor{lstString}{HTML}{008000}
\definecolor{lstNumber}{HTML}{0000FF}
\definecolor{lstComment}{HTML}{808080}
\definecolor{lstLineNumber}{HTML}{999999}
\lstdefinestyle{customhask}{ 
    belowcaptionskip=1\baselineskip, 
    breaklines=true, 
    frame=L, 
    xleftmargin=\parindent, 
    language=haskell, 
    backgroundcolor=\color{lstBG},
    commentstyle=\itshape\color{lstComment},
    keywordstyle=\bfseries\color{lstKeyword},
    numberstyle=\tiny\color{lstLineNumber},
    stringstyle=\color{lstString},
    basicstyle=\scriptsize\ttfamily,
    breakatwhitespace=false,        
    breaklines=false,               
    captionpos=b,                    
    keepspaces=true,              
    numbers=left,        
    numbersep=10pt,          
    showspaces=false,
    showstringspaces=false,
    showtabs=false, 
    tabsize=4,
    columns=flexible
} 
\lstdefinestyle{normalhask}{ 
    belowcaptionskip=1\baselineskip, 
    breaklines=true, 
    frame=L, 
    xleftmargin=\parindent, 
    language=haskell, 
    backgroundcolor=\color{lstBG},
    commentstyle=\itshape\color{lstComment},
    keywordstyle=\bfseries\color{lstKeyword},
    numberstyle=\tiny\color{lstLineNumber},
    stringstyle=\color{lstString},
    basicstyle=\normalsize\ttfamily,
    breakatwhitespace=false,        
    breaklines=false,               
    captionpos=b,                    
    keepspaces=true,              
    numbers=left,        
    numbersep=10pt,          
    showspaces=false,
    showstringspaces=false,
    showtabs=false, 
    tabsize=4,
    columns=flexible
} 
\lstset{escapechar=@,style=normalhask} 
\usepackage{xparse} 
\usepackage{xspace}
\newcommand{\includecodecap}[2][hask]{\lstinputlisting[caption=#2, escapechar=, style=custom#1]{#2}} 
\newcommand{\includecode}[2][hask]{\lstinputlisting[escapechar=@, style=custom#1]{#2}} 
\newcommand{\includelines}[4][hask]{\lstinputlisting[escapechar=, style=custom#1, firstline=#3, lastline=#4]{#2}}
\hypersetup{
    colorlinks=true,
    linkcolor=blue,
    filecolor=magenta,      
    urlcolor=blue,
    pdfpagemode=FullScreen,
    }

\DeclareMathOperator{\Term}{\mathcal{T}}
\DeclareMathOperator{\FlatTerm}{\mathcal{FT}}
\DeclareMathOperator{\Var}{\mathbf{Var}}
\DeclareMathOperator{\Cons}{\mathcal{C}}
\DeclareMathOperator{\Kan}{\mathcal{G}}
\DeclareMathOperator{\Fresh}{\mathbf{Fresh}}
\DeclareMathOperator{\Delay}{\mathbf{Delay}}
\DeclareMathOperator{\Call}{\mathbf{Call}}
\DeclareMathOperator{\Def}{\mathcal{D}}
\DeclareMathOperator{\Base}{\mathbf{Base}}
\DeclareMathOperator{\Conj}{\mathbf{Conj}}
\DeclareMathOperator{\free}{\mathbf{free}}
\DeclareMathOperator{\ground}{\mathbf{ground}}
\DeclareMathOperator{\In}{\mathbf{In}}
\DeclareMathOperator{\Out}{\mathbf{Out}}
\DeclareMathOperator{\mnondet}{\mathbf{nondet}}
\DeclareMathOperator{\mdet}{\mathbf{det}}
\DeclareMathOperator{\NotDet}{\mathbf{NotDet}}
\DeclareMathOperator{\BecameDet}{\mathbf{BecameDet}}
\DeclareMathOperator{\WasDet}{\mathbf{WasDet}}
\DeclareMathOperator{\Fun}{\mathcal{F}}
\DeclareMathOperator{\Return}{\mathbf{Return}}
\DeclareMathOperator{\Bind}{\mathbf{Bind}}
\DeclareMathOperator{\Match}{\mathbf{Match}}
\DeclareMathOperator{\Sum}{\mathbf{Sum}}
\DeclareMathOperator{\Guard}{\mathbf{Guard}}
\DeclareMathOperator{\Gen}{\mathbf{Gen}}
\DeclareMathOperator{\Stream}{\mathit{Stream}}
\DeclareMathOperator{\vars}{vars}
\DeclareMathOperator{\mode}{mode}
\DeclareMathOperator{\Mode}{\mathcal{M}}
\newcommand{\KanN}{\mathcal{K}^{N}}
\newcommand{\eps}{\varepsilon}
\newcommand{\TODO}[1][]{
  \vspace{0.2em}
  {\bf\color{red} TODO:} #1
  \vspace{0.2em}
}
\newcommand{\tran}[1]{\left\llbracket #1 \right\rrbracket}
\newcommand{\lst}[1]{\left[ #1 \right]}
\renewcommand{\emptyset}{\varnothing}